% Part: sets-functions-relations
% Chapter: sets
% Section: unions-and-intersections

\documentclass[../../../include/open-logic-section]{subfiles}

\begin{document}

\olfileid{sfr}{set}{uni}
\olsection{યોગગણો અને છેદગણો}

\begin{explain}
\olref[sfr][set][bas]{sec}માં આપણે ગુણધર્મ દ્વારા ગણોની વ્યાખ્યા,
એટલે કે $\Setabs{x}{\phi(x)}$ સ્વરૂપની વ્યાખ્યાઓ રજૂ કરી હતી.
અહીં આપણે કોઈ ગુણધર્મ~$\phi$નો ઉપયોગ કરીએ છીએ, અને આ ગુણધર્મમાં
અગાઉ વ્યાખ્યાયિત કરેલા ગણોનો ઉલ્લેખ હોઈ શકે છે.
દાખલા તરીકે, જો $A$ અને~$B$ ગણો હોય, તો
$\Setabs{x}{x \in A \lor x \in B}$ એ બધી વસ્તુઓનો ગણ છે જે $A$ અથવા~$B$ના
!!{element}s હોય; એટલે કે તે $A$ અને~$B$ના !!{element}sને એકત્ર કરતો ગણ છે.
આને \olref{fig:union}માં દર્શાવ્યા મુજબ જોઈ શકાય, જ્યાં રંગેલો વિસ્તાર
બંને ગણો $A$ અને~$B$ના !!{element}sને એકસાથે દર્શાવે છે.

\begin{figure}
  \olasset{assets/diagrams/union.tikz}
  \caption{બે ગણોનો યોગગણ $A \cup B$ એ $A$ના !!{element}s અને
  $B$ના !!{element}sને એકત્ર કરતાં મળતો ગણ છે.}
  \ollabel{fig:union}
\end{figure}

ગણો પરની આ ક્રિયા---તેમને એકત્ર કરવાની ક્રિયા---ખૂબ ઉપયોગી અને સામાન્ય છે,
તેથી આપણે તેને ઔપચારિક નામ અને સંકેત આપીએ છીએ.
\end{explain}

\begin{defn}[યોગગણ]
બે ગણો $A$ અને $B$નો \emph{યોગગણ}, જેને $A \cup B$ લખીએ છીએ,
એ બધી વસ્તુઓનો ગણ છે જે $A$ના, $B$ના, અથવા બંનેના !!{element}s હોય.
\[
A \cup B = \Setabs{x}{x \in A \lor x \in B}
\]
\end{defn}

\begin{ex}
!!{element}s કેટલી વાર આવ્યા છે તે મહત્ત્વનું ન હોવાથી, કોઈ !!a{element}
સામાન્ય હોય તેવા બે ગણોના યોગગણમાં તે !!{element} એક જ વાર આવે છે.
દા.ત., $\{ a, b, c\} \cup \{ a, 0, 1\} = \{a, b, c, 0, 1\}$.

ગણ અને તેના કોઈ ઉપગણનો યોગગણ મોટો ગણ પોતે જ છે:
$\{a, b, c \} \cup \{a \} = \{a, b, c\}$.

ગણ અને ખાલી ગણનો યોગગણ તે ગણ પોતે જ છે:
$\{a, b, c \} \cup \emptyset = \{a, b, c \}$.
\end{ex}

\begin{prob}
જો $A \subseteq B$ હોય, તો $A \cup B = B$ થાય તે સાબિત કરો.
\end{prob}

\begin{explain}
આપણે યોગક્રિયાની “દ્વૈત” ક્રિયાનો પણ વિચાર કરી શકીએ.
આ ક્રિયા એવા બધા !!{element}sનો ગણ બનાવે છે જે $A$ના !!{element}s
હોય અને $B$ના પણ !!{element}s હોય.
આ ક્રિયાને \emph{છેદક્રિયા} કહે છે અને તેને \olref{fig:intersection}માં
દર્શાવ્યા પ્રમાણે ચિત્રી શકાય.
\begin{figure}
  \olasset{assets/diagrams/intersection.tikz}
  \caption{બે ગણોનો છેદગણ $A \cap B$ એ તેમના સામાન્ય !!{element}sનો ગણ છે.}
  \ollabel{fig:intersection}
\end{figure}
\end{explain}

\begin{defn}[છેદગણ]
બે ગણો $A$ અને $B$નો \emph{છેદગણ}, જેને $A \cap B$ લખીએ છીએ,
એ બધી વસ્તુઓનો ગણ છે જે $A$ અને~$B$ બંનેના !!{element}s હોય.
\[
A \cap B = \Setabs{x}{x \in A \land x \in B}
\]
જો બે ગણોનો છેદગણ ખાલી હોય, તો તેમને \emph{અલગ ગણો} કહે છે.
આનો અર્થ એ છે કે તેમનામાં એક પણ સામાન્ય !!{element} નથી.
\end{defn}

\begin{ex}
જો બે ગણોમાં એક પણ સામાન્ય !!{element} ન હોય, તો તેમનો છેદગણ ખાલી છે:
$\{ a, b, c\} \cap \{ 0, 1\} = \emptyset$.

જો બે ગણોમાં સામાન્ય !!{element}s હોય, તો તેમનો છેદગણ એ બધા
!!{element}sનો ગણ છે: $\{a, b, c \} \cap \{a, b, d \} = \{a, b\}$.

ગણ અને તેના કોઈ ઉપગણનો છેદગણ નાનો ગણ પોતે જ છે:
$\{a, b, c\} \cap \{a, b\} = \{a, b\}$.

કોઈ પણ ગણ અને ખાલી ગણનો છેદગણ ખાલી છે:
$\{a, b, c \} \cap \emptyset = \emptyset$.
\end{ex}

\begin{prob}
જો $A \subseteq B$ હોય, તો $A \cap B = A$ થાય તે ચોકસાઈપૂર્વક સાબિત કરો.
\end{prob}

\begin{explain}
આપણે બે કરતાં વધુ ગણોનો યોગગણ કે છેદગણ પણ બનાવી શકીએ.
સામાન્ય રીતે આ કરવાની એક સુઘડ રીત નીચે મુજબ છે:
ધારો કે જે બધા ગણોનો યોગગણ (અથવા છેદગણ) બનાવવો છે તેમને
આપણે એક જ ગણમાં એકત્ર કરીએ. ત્યાર પછી, આપણા બધા મૂળ ગણોના
યોગગણની વ્યાખ્યા એ બધી વસ્તુઓના ગણ તરીકે કરી શકીએ જે આ ગણના
ઓછામાં ઓછા એક !!{element}માં હોય; અને છેદગણની વ્યાખ્યા એ બધી
વસ્તુઓના ગણ તરીકે કરી શકીએ જે આ ગણના દરેક !!{element}માં હોય.
\end{explain}

\begin{defn}
જો $A$ ગણોનો ગણ હોય, તો $\bigcup A$ એ $A$ના !!{element}sના
!!{element}sનો ગણ છે:
\begin{align*}
\bigcup A & = \Setabs{x}{x \text{ ગણ } A \text{ના કોઈ !!a{element}માં છે}},
\text{ એટલે કે,}\\
& = \Setabs{x}{\text{એવો } B \in A
  \text{ છે કે } x \in B}
\end{align*}
\end{defn}

\begin{defn}
જો $A$ ગણોનો ગણ હોય, તો $\bigcap A$ એ બધી વસ્તુઓનો ગણ છે જે
$A$ના બધા ઘટકોમાં સામાન્ય હોય:
\begin{align*}
\bigcap A & = \Setabs{x}{x \text{ ગણ } A \text{ના દરેક !!{element}માં છે}},
\text{ એટલે કે,}\\
 & = \Setabs{x}{\text{દરેક } B \in A, \text{ માટે } x \in B}
\end{align*}
\end{defn}

\begin{ex}
ધારો કે $A = \{ \{ a, b \}, \{ a, d, e \}, \{ a, d \} \}$.
તો $\bigcup A = \{ a, b, d, e \}$ અને $\bigcap A = \{ a \}$.
\end{ex}
\begin{prob}
જો $A$ ગણ હોય અને $A \in B$ હોય, તો $A \subseteq \bigcup B$ થાય તે બતાવો.
\end{prob}

ગણોના અનુક્રમ $A_1$, $A_2$, \dots માટે પણ આપણે આમ જ કરી શકીએ:
\begin{align*}
\bigcup_i A_i & = \Setabs{x}{x \text{ કોઈ એક } A_i \text{માં છે}}\\
\bigcap_i A_i & = \Setabs{x}{x \text{ દરેક } A_i \text{માં છે}}.
\end{align*}

જ્યારે આપણી પાસે ગણોનો \emph{સૂચકગણ} હોય, એટલે કે એવો ગણ $I$ હોય કે
દરેક $i \in I$ માટે આપણે $A_i$નો વિચાર કરતા હોઈએ,
ત્યારે આ સંક્ષિપ્ત સંજ્ઞાઓ પણ વાપરી શકીએ:
\begin{align*}
	\bigcup_{i \in I} A_i & = \bigcup \Setabs{A_i }{i \in I}\\
	\bigcap_{i \in I} A_i & = \bigcap\Setabs{A_i}{i \in I}
\end{align*}

છેલ્લે, આપણે $A$ના એવા બધા !!{element}sના ગણનો વિચાર કરવા માગીએ
જે $B$માં નથી. તેને \olref{difference}માં દર્શાવ્યા પ્રમાણે ચિત્રી શકાય.

\begin{figure}
  \olasset{assets/diagrams/difference.tikz}
  \caption{બે ગણોનો તફાવત ગણ $A \setminus B$ એ $A$ના એવા
  !!{element}sનો ગણ છે જે $B$ના પણ !!{element}s નથી.}
  \ollabel{difference}
\end{figure}

\begin{defn}[તફાવત ગણ]
\emph{તફાવત ગણ}~$A \setminus B$ એ $A$ના એવા બધા !!{element}sનો ગણ છે
જે $B$ના પણ !!{element}s નથી, એટલે કે,
\[
A\setminus B = \Setabs{x}{x\in A \text{ અને } x \notin B}.
\]
\end{defn}

\begin{prob}
જો $A \subsetneq B$ હોય, તો $B \setminus A \neq \emptyset$ થાય તે સાબિત કરો.
\end{prob}

\end{document}
